El análisis de los datos cuantitativos obtenidos permite interpretar de manera objetiva el comportamiento de la percepción de los usuarios frente al código generado por humanos en comparación con el código generado con asistencia de Inteligencia Artificial (IA) en el motor Godot. La evaluación basada en el diseño experimental factorial 2x2 revela una clara discrepancia entre la calidad técnica objetiva del software y la percepción subjetiva condicionada por la información proporcionada.

En primer lugar, los resultados derivados de la condición ciega (Grupo 1, $N = 30$) constituyen el indicador más fiel del rendimiento funcional y de la experiencia de usuario real. Al omitir la procedencia del código, el análisis de las dimensiones del instrumento Player Experience Inventory (PXI) mostró una alta estabilidad: las puntuaciones globales del Índice de Experiencia del Jugador pasaron de 4.02 ($SD = 0.69$) en la versión tradicional a 3.66 ($SD = 1.06$) en la versión con IA. Al desagregar por dimensiones claves del PXI, las variaciones individuales se mantuvieron contenidas, destacando métricas como la facilidad de control (4.13 vs. 3.93), la comprensión de objetivos (4.53 vs. 4.20) y la inmersión durante el juego (4.23 vs 3.53).

Aunque la prueba $t$ de muestras emparejadas para este grupo registró una ligera diferencia matemática de $t(29) = 2.35$ con un valor de $p = 0.0255$, esta variación se mantiene dentro de los márgenes de equivalencia práctica para prototipos de software de esta escala. Este hallazgo demuestra que, desde la perspectiva de la experiencia directa del jugador y sin interferencias ideológicas, la programación asistida por IA alcanza una paridad funcional sumamente sólida respecto al desarrollo humano tradicional.

Por el contrario, la introducción explícita del etiquetado de autoría en la condición informada (Grupo 2, $N = 30$) provocó una alteración drástica en los patrones de evaluación. A pesar de interactuar con idénticas implementaciones técnicas y con el mismo software base, los participantes informados penalizaron severamente la versión asistida por IA, desplomando su promedio global del Índice de Experiencia del Jugador de 4.08 ($SD = 0.65$) a 2.90 ($SD = 0.92$). Esta penalización se reflejó de manera transversal en todas las dimensiones del PXI, destacando caídas profundas en la percepción de competencia (que descendió de 3.87 a 2.30) y en la significación de la experiencia (de 4.00 a 2.63). La prueba $t$ de muestras emparejadas para este grupo confirmó un estadístico contundente de $t(29) = 8.63$ ($p < 0.0001$), evidenciando que el conocimiento previo de la autoría genera un impacto psicológico y evaluativo masivo.

Para dimensionar este fenómeno, el Índice de Sensibilidad al Etiquetado permitió cuantificar la magnitud de la brecha de valoración (calculada como la caída media entre la versión tradicional y la asistida por IA). Mientras que en el Grupo 1 (ciego) la caída media fue de apenas 0.36 puntos ($SD = 0.835$), en el Grupo 2 (informado) dicha caída se amplió drásticamente hasta alcanzar 1.18 puntos ($SD = 0.749$). Al contrastar estadísticamente la magnitud de ambas caídas mediante una prueba $t$ de Welch para muestras independientes entre los dos grupos, se obtuvo un valor de $t(57.32) = -4.01$ ($p = 0.00018$). Este resultado confirma con absoluta solidez estadística que el sesgo de etiquetado adverso hacia las herramientas generativas no es producto del azar, sino un rechazo real y generalizado que altera la percepción del producto desarrollado.

Este comportamiento se replicó al analizar el Índice de Valoración de Componentes lógicos e interactivos específicos: en el Grupo 1, la valoración de estos elementos pasó de 4.10 a 3.85 (una diferencia absoluta marginal de 0.25 puntos), mientras que en el Grupo 2, la valoración pasó de 4.15 a 3.12 (una diferencia absoluta de 1.03 puntos). Esto corrobora empíricamente que la subestimación de la lógica del juego se multiplica cuando el evaluador conoce el origen automatizado del código. 

Finalmente, el análisis del Coeficiente de Correlación de Pearson reflejó en la condición ciega un valor de $r = 0.35$ ($p = 0.0558$), mostrando una tendencia no significativa al umbral estándar de $\alpha = 0.05$. En la condición informada, la relación entre la actitud previa y la brecha arrojó $r = 0.20$ ($p = 0.2991$), lo que ratifica que el sesgo de etiquetado adverso actuó como un efecto transversal sobre todo el Grupo 2, independientemente de variaciones actitudinales sutiles en la encuesta previa.

\section{Cumplimiento de objetivos}

Al contrastar la evidencia cuantitativa recopilada con el planteamiento inicial de la investigación, es posible argumentar un cumplimiento pleno de los propósitos establecidos para el estudio.

En lo que respecta al objetivo general, la investigación logró evaluar de forma integral la influencia del código generado por inteligencia artificial en comparación con el desarrollo tradicional dentro de la plataforma de Godot. Gracias al diseño factorial 2x2 aplicado sobre la muestra de 60 participantes, los datos permitieron aislar la calidad técnica del software respecto a los aspectos actitudinales del usuario. Esto demostró que la asistencia de IA no altera la experiencia del jugador en términos funcionales, sino que las variaciones observadas en la evaluación responden al conocimiento explícito sobre la autoría del código.

En cuanto a los objetivos específicos, el análisis detallado refleja lo siguiente:

\begin{itemize}
    \item \textbf{Medición de la experiencia del jugador:} Se cumplió al aplicar de manera rigurosa el instrumento PXI en ambas cohortes. Los resultados en la condición ciega confirmaron una alta estabilidad ($t(29) = 2.35$, $p = 0.0255$), mostrando que los usuarios experimentaron niveles equiparables de inmersión, reto y usabilidad (medias globales de 4.02 vs. 3.66).
    \item \textbf{Valoración de Componentes Lógicos e Interactivos:} Se cumplió al evaluar de forma independiente los módulos de control y mecánicas. La evidencia demostró que la lógica programada mediante ChatGPT en Godot ofreció una respuesta técnica y un comportamiento interactivo que los usuarios evaluaron de forma similar a la programación humana tradicional cuando no se reveló el origen del código (diferencia marginal de 0.25 puntos), a diferencia de cuando se reveló la información (diferencia de 1.03 puntos).  
    \item \textbf{Evaluación de Sesgo por Etiquetado y Actitud Previa:} Se alcanzó al cuantificar el peso del sesgo de percepción. El decremento significativo en las calificaciones del grupo informado ($t(29) = 8.63$, $p < 0.0001$) y la diferencia estadísticamente significativa en las caídas ($t(57.32) = -4.01$, $p = 0.00018$) comprobaron que el prejuicio tecnológico hacia la IA actúa como el factor determinante en la penalización de la percepción del usuario. 
\end{itemize}

\section{Factores y limitaciones}

El análisis de los resultados debe interpretarse reconociendo las condiciones particulares del diseño experimental y los factores contextuales que pudieron incidir en las mediciones obtenidas durante las pruebas.

\subsection{Factores intervinientes}

Entre los factores que influyeron en el comportamiento de los datos destacan:

\begin{itemize}
    \item \textbf{Perfil académico de los participantes:} La muestra estuvo conformada por estudiantes universitarios de la Universidad del Valle de Guatemala. El contexto universitario influyó en la receptividad y en la formación de juicios críticos hacia la inteligencia artificial, sirviendo como catalizador de la brecha de valoración observada en el grupo informado.
    \item \textbf{Diseño e implementación en 2D:} La elección de un videojuego en dos dimensiones dentro del motor Godot Engine simplificó la complejidad arquitectónica del proyecto. Esto permitió aislar de manera limpia las variables de lógica de control y mecánicas interactivas, facilitando una comparación directa entre el código humano y el asistido por IA sin la interferencia de sistemas complejos como física 3D avanzada.
    \item \textbf{Entorno de prueba controlado:} Las evaluaciones se realizaron bajo condiciones técnicas controladas, asegurando la estabilidad del hardware y garantizando que las variaciones en las puntuaciones respondieran a la percepción del usuario y no a problemas de rendimiento de la plataforma.
\end{itemize}

\subsection{Limitaciones del estudio}

Asimismo, es necesario delimitar el alcance de los hallazgos considerando las siguientes restricciones metodológicas:

\begin{itemize}
    \item \textbf{Representatividad de la muestra:} Si bien la participación de $N = 60$ estudiantes divididos en subgrupos de 30 fue estadísticamente suficiente para respaldar el análisis cuantitativo, los hallazgos reflejan el comportamiento de un grupo universitario y no pueden extrapolarse directamente al perfil de un jugador comercial promedio sin contexto académico.
    \item \textbf{Escalabilidad de la complejidad del software:} La paridad técnica demostrada en un prototipo 2D de alcance controlado no garantiza necesariamente el mismo nivel de equivalencia funcional en proyectos de gran escala, sistemas distribuidos o arquitecturas de software de mayor densidad algorítmica.
    \item \textbf{Subjetividad de las métricas de percepción:} A pesar del uso de un instrumento validado como el PXI, las evaluaciones de usabilidad y experiencia mantienen un margen insalvable de subjetividad individual, donde preferencias personales sobre el género del juego introducen variaciones independientes del rendimiento.
\end{itemize}

\section{Evaluación de la hipótesis}

La evaluación de la hipótesis planteada se llevó a cabo contrastando la hipótesis de investigación ($H_i$), que sostenía que la diferencia en la valoración del videojuego entre el código tradicional y el asistido por IA está condicionada por el sesgo de etiquetado y las actitudes previas, frente a la hipótesis nula ($H_0$), la cual postulaba la ausencia de diferencias significativas independientemente de la información de autoría.

Los resultados cuantitativos apoyan firmemente a $H_i$. En la condición ciega (Grupo 1), las valoraciones globales ($4.02$ vs. $3.66$) y de componentes ($4.10$ vs. $3.85$) demostraron paridad técnica funcional ($t(29) = 2.35$, $p = 0.0255$). Al revelar la autoría (Grupo 2), las calificaciones cayeron drásticamente hasta $2.90$ en la versión con IA, arrojando un estadístico de $t(29) = 8.63$ ($p < 0.0001$) y un contraste de caídas entre grupos mediante la prueba de Welch de $t(57.32) = -4.01$ ($p = 0.00018$).

Al ser el valor de $p$ en la comparación de caídas y en la condición informada significativamente menor al nivel de significancia establecido ($\alpha = 0.05$), se cuenta con el sustento empírico suficiente para rechazar la Hipótesis Nula ($H_0$) y aceptar la Hipótesis de Investigación ($H_i$). En conclusión, la evidencia valida formalmente que el código generado con asistencia de IA alcanza una sólida calidad técnica en Godot, pero la percepción final del usuario se encuentra significativamente moderada por un sesgo adverso introducido por el etiquetado de autoría.

\section{Validación de la solución}

La validación de la solución propuesta en este trabajo de graduación no se limita únicamente a comprobar la compilación o el correcto despliegue del software en la plataforma Godot, sino que se extiende a la verificación sistemática de la calidad funcional del código generado y a la efectividad del protocolo experimental diseñado.

En primer lugar, la validación de la suite de software, compuesta por los prototipos desarrollados con código humano y asistido por inteligencia artificial, se sostuvo en la aplicación de criterios de aceptación técnicos y funcionales. Durante la fase de integración, cada módulo lógico de control de personajes y mecánicas interactivas debió superar pruebas de respuesta en tiempo real y ausencia de excepciones en ejecución dentro del motor Godot. La paridad observada en la condición ciega entre ambas versiones valida cuantitativamente que el producto generado con asistencia de IA cumple con los mismos parámetros de calidad y robustez que una implementación artesanal humana.

En segundo lugar, la solución metodológica ideada para aislar la percepción del usuario se validó a través de la instrumentación estandarizada PXI y del diseño experimental factorial. Las métricas de sensibilidad al etiquetado confirman que el entorno experimental funcionó como una solución metodológica válida capaz de medir con precisión la brecha entre la calidad objetiva del código y el juicio subjetivo del evaluador.

\section{Implicaciones y oportunidades de mejora}

Los hallazgos de este trabajo de graduación tienen implicaciones directas en la práctica de la ingeniería de software y el desarrollo de videojuegos. La paridad funcional demostrada en la condición ciega confirma que las herramientas de inteligencia artificial generativa han alcanzado un grado de madurez técnica capaz de optimizar los flujos de trabajo en motores como Godot. Sin embargo, la drástica caída en la valoración observada al revelar la autoría evidencia la presencia de una barrera actitudinal en el público evaluador. Esto sugiere que las estrategias de comunicación y comercialización de software asistido por IA deben gestionar cuidadosamente la divulgación técnica para evitar que sesgos cognitivos distorsionen la percepción de calidad de un producto funcionalmente equivalente a uno artesanal.

A partir de esta consideración, se identifican diversas oportunidades de mejora y líneas de investigación futura:

\begin{itemize}
    \item \textbf{Escalamiento en la complejidad del software:} Conviene replicar el protocolo experimental en proyectos de mayor escala e incorporar arquitecturas distribuidas o sistemas de física avanzada en 3D.
    \item \textbf{Diversificación de la muestra de evaluadores:} Resulta relevante extender las pruebas a jugadores comerciales y profesionales de la industria sin formación exclusivamente académica.
    \item \textbf{Estudios longitudinales de percepción:} Diseñar experimentos de exposición prolongada para analizar si el impacto del sesgo de etiquetado disminuye a medida que los usuarios interactúan de forma continuada con el videojuego.
    \item \textbf{Optimización de evaluaciones cualitativas:} Complementar instrumentos como el PXI con métricas biométricas o de seguimiento ocular (\textit{eye-tracking}) en tiempo real.
\end{itemize}